% !TEX root = ../../ctfp-print.tex

\lettrine[lhang=0.17]{圏}{論では、すべてのものが}すべてのものと関係していて、
あらゆるものが多くの角度から見られるように思える。例えば、
\hyperref[products-and-coproducts]{積}の普遍的構成を取り上げてみよう。
\hyperref[functors]{関手}と\hyperref[natural-transformations]{自然変換}について
さらに知った今、それを単純化し、できれば一般化できるだろうか？試してみよう。

\begin{figure}[H]
  \centering
  \includegraphics[width=0.3\textwidth]{images/productpattern.jpg}
\end{figure}

\noindent
積の構成は、積を構成したい二つの対象 $a$ と $b$ の選択から始まる。
しかし、\emph{対象を選択する}とはどういう意味か？この操作を、より圏論的な言葉で言い換えられるだろうか？
二つの対象はパターンを形成する――非常に単純なパターンだ。このパターンを圏として抽象化できる――
非常に単純な圏だが、それでも圏だ。これを $\cat{2}$ と呼ぼう。この圏には二つの対象 $1$ と $2$ だけがあり、
二つの必須な恒等射以外の射はない。これで、$\cat{C}$ における二つの対象の選択を、
圏 $\cat{2}$ から $\cat{C}$ への関手 $D$ を定義する行為として言い換えられる。
関手は対象を対象に写すので、その像はちょうど二つの対象だ（関手が対象を潰すなら一つになることもあるが、
それも問題ない）。また射も写す――この場合、単純に恒等射を恒等射に写す。

このアプローチの優れた点は、圏論的な概念を基に構築されていて、祖先の狩猟採集者の語彙から直接取り出したような
「対象を選択する」といった不正確な記述を避けている点だ。また、偶然にも一般化も容易だ。
なぜなら $\cat{2}$ より複雑な圏を使ってパターンを定義することを妨げるものは何もないからだ。

\begin{figure}[H]
  \centering
  \includegraphics[width=0.35\textwidth]{images/two.jpg}
\end{figure}

\noindent
しかし続けよう。積の定義の次のステップは、候補対象 $c$ の選択だ。
ここでも、単元圏からの関手の観点から選択を言い換えられる。実際、Kan拡張を使うなら、
それが正しい方法だろう。しかし、まだKan拡張の準備ができていないので、別のトリックが使える：
同じ圏 $\cat{2}$ から $\cat{C}$ への定数関手 $\Delta$ だ。$\cat{C}$ における $c$ の選択は
$\Delta_c$ で行える。$\Delta_c$ はすべての対象を $c$ に写し、すべての射を $\idarrow[c]$ に写すことを思い出そう。

\begin{figure}[H]
  \centering
  \includegraphics[width=0.35\textwidth]{images/twodelta.jpg}
\end{figure}

\noindent
今、$\cat{2}$ と $\cat{C}$ の間を行く二つの関手 $\Delta_c$ と $D$ がある。
だからそれらの間の自然変換を問うのは自然なことだ。$\cat{2}$ には対象が二つしかないので、
自然変換は二つの成分を持つ。$\cat{2}$ の対象 $1$ は $\Delta_c$ によって $c$ に写され、
$D$ によって $a$ に写される。だから $\Delta_c$ と $D$ の間の自然変換の $1$ における成分は、
$c$ から $a$ への射だ。これを $p$ と呼ぼう。同様に、二つ目の成分は $c$ から $b$ への射 $q$ だ――
$\cat{2}$ の対象 $2$ の $D$ による像だ。しかしこれらはまさに積の元の定義で使った
二つの射影と同じものだ。だから対象と射影の選択について語る代わりに、
関手と自然変換を選ぶことについて語るだけでよい。この単純なケースでは、
$\cat{2}$ には（恒等射以外の）射がないので、変換の自然性条件が自明に満たされることになる。

\begin{figure}[H]
  \centering
  \includegraphics[width=0.35\textwidth]{images/productcone.jpg}
\end{figure}

\noindent
この構成を $\cat{2}$ 以外の圏――例えば自明でない射を含むもの――へ一般化すると、
$\Delta_c$ と $D$ の間の変換に自然性条件が課される。
このような変換を\emph{錐}と呼ぶ。$\Delta$ の像が錐・ピラミッドの頂点であり、
その側面は自然変換の成分によって形成されるからだ。$D$ の像が錐の底面を形成する。

一般に、錐を構成するには、パターンを定義する圏 $\cat{I}$ から始める。
これは小さな圏で、しばしば有限圏だ。$\cat{I}$ から $\cat{C}$ への関手 $D$ を選んで、
それ（またはその像）を\emph{図式}と呼ぶ。錐の頂点として $\cat{C}$ の何らかの $c$ を選ぶ。
それを使って $\cat{I}$ から $\cat{C}$ への定数関手 $\Delta_c$ を定義する。
$\Delta_c$ から $D$ への自然変換が錐だ。有限の $\cat{I}$ では、これは単に
$c$ を図式：$D$ の下の $\cat{I}$ の像に結ぶ射の束だ。

\begin{figure}[H]
  \centering
  \includegraphics[width=0.35\textwidth]{images/cone.jpg}
\end{figure}

\noindent
自然性は、この図式のすべての三角形（ピラミッドの壁）が可換であることを要求する。
実際、$\cat{I}$ の任意の射 $f$ を取ろう。関手 $D$ はそれを $\cat{C}$ の射 $D f$ に写す。
この射がある三角形の底辺を形成する。定数関手 $\Delta_c$ は $f$ を $c$ 上の恒等射に写す。
$\Delta$ は射の両端を一つの対象に潰し、自然性の四角形が可換な三角形になる。
この三角形の二つの腕が自然変換の成分だ。

\begin{figure}[H]
  \centering
  \includegraphics[width=0.35\textwidth]{images/conenaturality.jpg}
\end{figure}

\noindent
これが一つの錐だ。我々が興味を持つのは\newterm{普遍錐}だ――積の定義のために普遍的な対象を選んだのと同様に。

これにはさまざまなアプローチがある。例えば、与えられた関手 $D$ に基づいた\emph{錐の圏}を定義できる。
その圏の対象は錐だ。ただし、$\cat{C}$ のすべての対象 $c$ が錐の頂点になれるわけではない。
なぜなら $\Delta_c$ と $D$ の間に自然変換が存在しないこともあるからだ。

これを圏にするには、錐の間の射も定義しなければならない。それらは頂点の間の射によって完全に決まる。
しかし、どんな射でもよいわけではない。積の構成において、候補対象（頂点）の間の射は
射影の共通因子でなければならないという条件を課したことを思い出そう。例えば：

\src{snippet01}

\begin{figure}[H]
  \centering
  \includegraphics[width=0.35\textwidth]{images/productranking.jpg}
\end{figure}

この条件は、一般の場合、一辺が因子分解射である三角形がすべて可換であるという条件に翻訳される。

\begin{figure}[H]
  \centering
  \includegraphics[width=0.35\textwidth]{images/conecommutativity.jpg}
  \caption{二つの錐を結ぶ可換三角形。因子分解射 $h$ を伴う
    （ここで、下の錐が $\Lim[D]$ を頂点とする普遍錐）}
\end{figure}

\noindent
これらの因子分解射を錐の圏における射として採用する。これらの射が確かに合成可能であること、
また恒等射も因子分解射であることを確認するのは容易だ。したがって錐は圏を形成する。

これで普遍錐を錐の圏における\emph{終対象}として定義できる。終対象の定義は、
他の任意の対象からその対象への唯一の射が存在すると述べている。この場合、
それは他の任意の錐の頂点から普遍錐の頂点への唯一の因子分解射が存在することを意味する。
この普遍錐を図式 $D$ の\emph{極限}、$\Lim[D]$ と呼ぶ（文献では、$\Lim$ の下に $I$ を指す
左向き矢印をよく見る）。しばしば略記として、この錐の頂点を極限（または極限対象）と呼ぶ。

直観としては、極限は図式全体の性質を一つの対象に体現する。例えば、二対象図式の極限が
二つの対象の積だ。積は（二つの射影とともに）両方の対象についての情報を含む。
また普遍的であるということは、余分なものを持たないことを意味する。

\section{自然同型としての極限}

この極限の定義にはまだ何か不満な点がある。つまり、機能するには機能するが、
任意の二つの錐を結ぶ三角形の可換性条件がまだ残っている。何らかの自然性条件に
置き換えられれば、はるかに優雅だろう。しかしどうすれば？

今、一つの錐だけでなく、錐の集合（実際には圏）全体を扱っている。極限が存在するなら
（そして――明確にしておくが――その保証はない）、それらの錐の一つが普遍錐だ。
他の任意の錐に対して、その頂点（$c$ と呼ぼう）を普遍錐の頂点（$\Lim[D]$ と名付けた）に
写す唯一の因子分解射がある。（実際、「他の」という言葉は省略できる。
なぜなら恒等射が普遍錐をそれ自身に写し、自明に自身を通じて因子分解されるからだ。）
重要な部分を繰り返そう：任意の錐が与えられると、特殊な種類の唯一の射がある。
錐から特殊な射への写像があり、それは一対一の写像だ。

この特殊な射はホム集合 $\cat{C}(c, \Lim[D])$ の元だ。このホム集合の他の元は
恵まれていない。つまり、二つの錐の写像を因子分解しないという意味で。
欲しいのは、各 $c$ に対して、集合 $\cat{C}(c, \Lim[D])$ から
特定の可換性条件を満たす一つの射を選ぶ能力だ。それは自然変換を定義するように聞こえないだろうか？
まさにそう聞こえる！

しかし、この変換によって関係付けられる関手は何か？

一方の関手は $c$ を集合 $\cat{C}(c, \Lim[D])$ に写す写像だ。これは $\cat{C}$ から $\Set$ への関手だ――
対象を集合に写す。実際、これは反変関手だ。射に対するその作用を定義しよう：
$c'$ から $c$ への射 $f$ を取ろう：
\[f \Colon c' \to c\]
この関手は $c'$ を集合 $\cat{C}(c', \Lim[D])$ に写す。$f$ に対するこの関手の作用を定義する
（言い換えれば、$f$ を持ち上げる）ために、$\cat{C}(c, \Lim[D])$ と
$\cat{C}(c', \Lim[D])$ の間の対応する写像を定義しなければならない。
では $\cat{C}(c, \Lim[D])$ の元 $u$ を一つ取って、それを
$\cat{C}(c', \Lim[D])$ の元に写せるか見てみよう。ホム集合の元は射なので：
\[u \Colon c \to \Lim[D]\]
$u$ に $f$ を前合成して得られる：
\[u \circ f \Colon c' \to \Lim[D]\]
これが $\cat{C}(c', \Lim[D])$ の元だ――射の写像が実際に見つかった：

\src{snippet02}
\emph{反変}関手に特有の $c$ と $c'$ の順序の逆転に注意せよ。

\begin{figure}[H]
  \centering
  \includegraphics[width=0.35\textwidth]{images/homsetmapping.jpg}
\end{figure}

\noindent
自然変換を定義するには、やはり $\cat{C}$ から $\Set$ への写像である別の関手が必要だ。
しかし今度は錐の集合を考える。錐は自然変換に過ぎないので、
自然変換の集合 $\mathit{Nat}(\Delta_c, D)$ を見ている。$c$ からこの特定の自然変換の集合への
写像は（反変）関手だ。どうやってそれを示すか？再び、射に対するその作用を定義しよう：
\[f \Colon c' \to c\]
$f$ の持ち上げは、$\cat{I}$ から $\cat{C}$ へ行く二つの関手の間の自然変換の写像でなければならない：
\[\mathit{Nat}(\Delta_c, D) \to \mathit{Nat}(\Delta_{c'}, D)\]
自然変換をどう写すか？各自然変換は射の選択だ――その成分――$\cat{I}$ の各元に一つの射。
$\mathit{Nat}(\Delta_c, D)$ の元である $\alpha$ の $a$（$\cat{I}$ の対象）における成分は射：
\[\alpha_a \Colon \Delta_c a \to D a\]
つまり、定数関手 $\Delta$ の定義を使えば、
\[\alpha_a \Colon c \to D a\]
$f$ と $\alpha$ が与えられたとき、$\mathit{Nat}(\Delta_{c'}, D)$ の元である $\beta$ を構成しなければならない。
$a$ における成分は射でなければならない：
\[\beta_a \Colon c' \to D a\]
前者（$\alpha_a$）に $f$ を前合成することで後者（$\beta_a$）を簡単に得られる：
\[\beta_a = \alpha_a \circ f\]
これらの成分が確かに自然変換をなすことを示すのは比較的容易だ。

\begin{figure}[H]
  \centering
  \includegraphics[width=0.4\textwidth]{images/natmapping.jpg}
\end{figure}

\noindent
射 $f$ が与えられたことで、成分ごとに二つの自然変換の間の写像を構築した。
この写像は関手の \code{contramap} を定義する：
\[c \to \mathit{Nat}(\Delta_c, D)\]
今やったことは、$\cat{C}$ から $\Set$ への（反変）関手が二つあることを示したことだ。
何も仮定していない――これらの関手は常に存在する。

ちなみに、これらの関手の最初のものは圏論において重要な役割を果たし、
米田の補題について話すときにまた見ることになる。任意の圏 $\cat{C}$ から $\Set$ への反変関手には
名前がある：「前層」と呼ばれる。これは\newterm{表現可能な前層}と呼ばれる。
二つ目の関手も前層だ。

二つの関手があるので、それらの間の自然変換について語れる。
では早速、結論を述べよう：$\cat{I}$ から $\cat{C}$ への関手 $D$ が極限 $\Lim[D]$ を持つのは、
今定義した二つの関手の間に自然同型が存在する場合、かつその場合に限る：
\[\cat{C}(c, \Lim[D]) \simeq \mathit{Nat}(\Delta_c, D)\]
自然同型とは何かを思い出そう。それはすべての成分が同型射、すなわち可逆な射である自然変換だ。

この命題の証明には立ち入らない。手順は退屈ではあるが、かなり直截的だ。
自然変換を扱うとき、通常は成分、つまり射に焦点を当てる。この場合、両方の関手の値域が $\Set$ なので、
自然同型の成分は関数になる。これらは高階関数だ。ホム集合から自然変換の集合へ行くからだ。
再び、関数が引数に何をするかを考えることで分析できる：ここで引数は射――
$\cat{C}(c, \Lim[D])$ の元――であり、結果は自然変換――$\mathit{Nat}(\Delta_c, D)$ の元、
つまり錐と呼んできたもの――だ。この自然変換は、射である独自の成分を持つ。
つまり結局のところ射ばかりで、それらを追跡できれば命題を証明できる。

最も重要な結果は、この同型の自然性条件が、錐の写像の可換性条件とまったく同じであることだ。

今後の展開の予告として、$\mathit{Nat}(\Delta_c, D)$ の集合が関手圏のホム集合として考えられることに
触れておこう。だから自然同型は二つのホム集合を関係付けるが、それはアジャンクションと呼ばれる
さらに一般的な関係を指し示す。

\section{極限の例}

$\cat{2}$ と呼んだ単純な圏によって生成される図式の極限が圏論的積であることを見た。

さらに単純な極限の例として終対象がある。最初は単元圏が終対象につながると考えがちだが、
実情はさらにシンプルだ：終対象は空圏によって生成される極限だ。
空圏からの関手は対象を選ばないので、錐は頂点だけに縮小する。
普遍錐は、他の任意の頂点からそこへの唯一の射を持つ唯一の頂点だ。
これが終対象の定義だと気づくだろう。

次の興味深い極限は\emph{等化子}と呼ばれる。これは二つの並行した射がそれらの間を行く
二要素圏（そしていつものように恒等射も）によって生成される極限だ。
この圏は $\cat{C}$ において、二つの対象 $a$ と $b$、および二つの射からなる図式を選ぶ：

\src{snippet03}

この図式の上に錐を構成するには、頂点 $c$ と二つの射影を追加しなければならない：

\src{snippet04}

\begin{figure}[H]
  \centering
  \includegraphics[width=0.35\textwidth]{images/equalizercone.jpg}
\end{figure}

\noindent
可換でなければならない二つの三角形がある：

\src{snippet05}

これは $q$ がこれらの等式の一つ、例えば \code{q = f . p} によって一意に決まることを示しており、
図から省略できる。これで条件が一つだけ残る：

\src{snippet06}

$\Set$ に限定して考えると、関数 $p$ の像は $a$ の部分集合を選ぶ。この部分集合に制限すると、
関数 $f$ と $g$ は等しい。

例えば、$a$ を座標 $x$ と $y$ でパラメータ化された二次元平面とする。
$b$ を実数直線とし、次のように取ろう：

\src{snippet07}

これら二つの関数の等化子は実数の集合（頂点 $c$）と関数：

\src{snippet08}

だ。$(p~t)$ は二次元平面上の直線を定義することに注意せよ。この直線に沿って、
二つの関数は等しい。

もちろん、等式につながりうる他の集合 $c'$ と関数 $p'$ もある：

\src{snippet09}

しかしそれらはすべて $p$ を通じて一意に因子分解される。例えば、単元集合 $\cat{()}$ を
$c'$ として取り、関数：

\src{snippet10}

とすることができる。$f (0, 0) = g (0, 0)$ なので、これは良い錐だ。
しかし普遍的ではない。なぜなら $h$ を通じた唯一の因子分解があるからだ：

\src{snippet11}

ただし

\src{snippet12}

\begin{figure}[H]
  \centering
  \includegraphics[width=0.35\textwidth]{images/equilizerlimit.jpg}
\end{figure}

\noindent
こうして等化子は $f~x = g~x$ 型の方程式を解くために使える。しかしそれはより一般的だ。
なぜなら代数的にではなく、対象と射の観点から定義されているからだ。

方程式を解くというさらに一般的なアイデアは、別の極限――引き戻し――に体現されている。
ここでも等しくしたい二つの射があるが、今度はそれらの定義域が異なる。
形が $1\rightarrow2\leftarrow3$ の三対象圏から始める。この圏に対応する図式は
三つの対象 $a$、$b$、$c$ と二つの射からなる：

\src{snippet13}

この図式は\emph{コスパン}と呼ばれることが多い。

この図式の上に構成される錐は頂点 $d$ と三つの射からなる：

\src{snippet14}

\begin{figure}[H]
  \centering
  \includegraphics[width=0.35\textwidth]{images/pullbackcone.jpg}
\end{figure}

\noindent
可換性条件は $r$ が他の射によって完全に決まることを示しており、図から省略できる。
だから次の条件だけが残る：

\src{snippet15}
引き戻しはこの形の普遍錐だ。

\begin{figure}[H]
  \centering
  \includegraphics[width=0.35\textwidth]{images/pullbacklimit.jpg}
\end{figure}

\noindent
再び集合に限定して考えると、対象 $d$ を $a$ と $c$ の要素のペアであって、
最初の成分に作用する $f$ が第二の成分に作用する $g$ と等しいものからなると考えられる。
それでもまだ抽象的すぎるなら、$g$ が定数関数、例えば $g~\_ = 1.23$ の特殊ケースを考えよう
（$b$ が実数の集合と仮定して）。すると方程式を実際に解くことになる：

\src{snippet16}

この場合、$c$ の選択は（空集合でなければ）無関係なので、単元集合として取ることができる。
例えば $a$ を三次元ベクトルの集合、$f$ をベクトルの長さとする。すると引き戻しは
ペア $(v, ())$ の集合になる。ここで $v$ は長さ 1.23 のベクトル
（方程式 $\sqrt{(x^{2}+y^{2}+z^{2})} = 1.23$ の解）であり、
$()$ は単元集合のダミー要素だ。

しかし引き戻しにはプログラミングでもより一般的な応用がある。例えば、
射がサブクラスからスーパークラスへ結ぶ矢印である C++ クラスの圏を考えよう。
継承を推移的な性質と考える。\code{C} が \code{B} を継承し、\code{B} が \code{A} を継承するなら
\code{C} は \code{A} を継承すると言う（結局、\code{A} へのポインタが期待される場所に
\code{C} へのポインタを渡せる）。また \code{C} は \code{C} を継承すると仮定するので、
すべてのクラスに恒等矢印がある。これによりサブクラス化はサブタイピングと一致する。
C++ は多重継承もサポートするので、\code{B} と \code{C} が \code{A} を継承し、
4番目のクラス \code{D} が \code{B} と \code{C} を多重継承するダイヤモンド継承図式を構成できる。
通常、\code{D} は \code{A} の二つのコピーを持つことになるが、これはめったに望ましくない。
しかし仮想継承を使えば \code{D} の中に \code{A} のコピーを一つだけ持てる。

この図式で \code{D} が引き戻しであるとはどういうことか？それは \code{B} と \code{C} を多重継承する
任意のクラス \code{E} が \code{D} のサブクラスでもあることを意味する。これは C++ では直接表現できない。
C++ ではサブタイピングが名目的だからだ（C++ のコンパイラはこの種のクラス関係を推論しない――
「ダックタイピング」が必要になる）。しかしサブタイピングの関係の外に出て、代わりに
\code{E} から \code{D} へのキャストが安全かどうかを問うことができる。
このキャストは \code{D} が \code{B} と \code{C} の最小限の組み合わせで、
追加データなし、メソッドのオーバーライドなしなら安全だろう。もちろん、
\code{B} と \code{C} の一部のメソッド間に名前の衝突があれば引き戻しは存在しない。

\begin{figure}[H]
  \centering
  \includegraphics[width=0.25\textwidth]{images/classes.jpg}
\end{figure}

\noindent
型推論においても引き戻しのより高度な使用法がある。二つの式の型を\emph{単一化}する必要が
しばしばある。例えば、コンパイラが関数の型を推論したいとする：

\begin{snip}{haskell}
twice f x = f (f x)
\end{snip}
コンパイラはすべての変数と部分式に仮の型を割り当てる。特に：

\begin{snip}{haskell}
f       :: t0
x       :: t1
f x     :: t2
f (f x) :: t3
\end{snip}
そこから導出する：

\begin{snip}{haskell}
twice :: t0 -> t1 -> t3
\end{snip}
また、関数適用の規則から生じる制約の集合も出てくる：

\begin{snip}{haskell}
t0 = t1 -> t2 -- f が x に適用されるから
t0 = t2 -> t3 -- f が (f x) に適用されるから
\end{snip}
これらの制約は、両方の式の未知の型に代入したとき同じ型を生成する型（または型変数）の
集合を見つけることで単一化しなければならない。そのような代入の一つは：

\begin{snip}{haskell}
t1 = t2 = t3 = Int
twice :: (Int -> Int) -> Int -> Int
\end{snip}
しかし明らかに、これは最も一般的なものではない。最も一般的な代入は引き戻しを使って得られる。
詳細には立ち入らない――本書の範囲を超えているから――しかし結果は次のようになることを
自分で確かめられる：

\begin{snip}{haskell}
twice :: (t -> t) -> t -> t
\end{snip}
ここで \code{t} は自由型変数だ。

\section{余極限}

圏論のすべての構成と同様、極限にも反対圏での双対像がある。錐のすべての矢印の向きを逆にすると、
余錐が得られ、その普遍的なものを余極限と呼ぶ。逆転は因子分解射にも影響することに注意せよ。
今度は普遍余錐から他の余錐へ流れる。

\begin{figure}[H]
  \centering
  \includegraphics[width=0.35\textwidth]{images/colimit.jpg}
  \caption{二つの頂点を結ぶ因子分解射 $h$ を持つ余錐。}
\end{figure}

\noindent
余極限の典型的な例は余積だ。積の定義で使った圏 $\cat{2}$ によって生成される図式に対応する。

\begin{figure}[H]
  \centering
  \includegraphics[width=0.35\textwidth]{images/coproductranking.jpg}
\end{figure}

\noindent
積も余積も、それぞれ異なる方法で対象のペアの本質を体現する。

終対象が極限であったように、始対象は空圏に基づく図式に対応する余極限だ。

引き戻しの双対は\emph{押し出し}と呼ばれる。これは圏
$1\leftarrow2\rightarrow3$ によって生成されるスパンと呼ばれる図式に基づいている。

\section{連続性}

以前、関手は既存の接続（射）を決して壊さないという意味で、圏の連続写像のアイデアに
近いと述べた。圏 $\cat{C}$ から $\cat{C'}$ への\emph{連続関手} $F$ の実際の定義には、
関手が極限を保存するという要件が含まれる。$\cat{C}$ の任意の図式 $D$ は、
単純に二つの関手を合成することで $\cat{C'}$ の図式 $F \circ D$ に写せる。
$F$ の連続性条件は、図式 $D$ が極限 $\Lim[D]$ を持つなら、
図式 $F \circ D$ も極限を持ち、それが $F (\Lim[D])$ に等しいと述べる。

\begin{figure}[H]
  \centering
  \includegraphics[width=0.6\textwidth]{images/continuity.jpg}
\end{figure}

\noindent
関手は射を射に写し、合成を合成に写すので、錐の像は常に錐であることに注意せよ。
可換三角形は常に可換三角形に写される（関手は合成を保存する）。
同じことが因子分解射にも当てはまる：因子分解射の像もまた因子分解射だ。
だからすべての関手は\emph{ほとんど}連続だ。うまくいかないかもしれないのは唯一性条件だ。
$\cat{C'}$ の因子分解射は一意でないかもしれない。また $\cat{C'}$ には $\cat{C}$ では
利用できなかった「より良い錐」が他にある可能性もある。

ホム関手は連続関手の例だ。ホム関手 $\cat{C}(a, b)$ は最初の変数で反変で、
二番目の変数で共変であることを思い出そう。言い換えれば、それは関手：
\[\cat{C}^\mathit{op} \times \cat{C} \to \Set\]
だ。二番目の引数が固定されると、ホム集合関手（表現可能な前層になる）は
$\cat{C}$ の余極限を $\Set$ の極限に写す。最初の引数が固定されると、極限を極限に写す。

Haskell では、ホム関手は任意の二つの型を関数型に写す写像なので、
単にパラメータ化された関数型だ。二番目のパラメータを、例えば \code{String} に固定すると、
反変関手が得られる：

\src{snippet17}
連続性は、\code{ToString} が余極限、例えば余積 \code{Either b c} に適用されると、
極限を生成することを意味する。この場合、二つの関数型の積：

\src{snippet18}
確かに、\code{Either b c} の任意の関数は、二つのケースがそれぞれ関数のペアによって
処理される場合分け節として実装される。

同様に、ホム集合の最初の引数を固定すると、おなじみのリーダー関手が得られる。
その連続性は、例えば、積を返す任意の関数が関数の積と同値であることを意味する。特に：

\src{snippet19}
あなたが何を考えているか分かる：これらのことを理解するのに圏論は必要ない。
そうかもしれない！それでも、このような結果がビットとバイト、プロセッサアーキテクチャ、
コンパイラ技術、さらにはラムダ計算にも頼らずに、第一原理から導出できることには驚嘆する。

「極限」と「連続性」という名前がどこから来たか知りたければ、それらは微積分の対応する概念の
一般化だ。微積分では、極限と連続性は開近傍の観点から定義される。
位相を定義する開集合は圏（半順序集合）を形成する。

\section{課題}

\begin{enumerate}
  \tightlist
  \item
        C++ クラスの圏で押し出しをどう記述するか？
  \item
        恒等関手 $\mathbf{Id} \Colon \cat{C} \to \cat{C}$ の極限が始対象であることを示せ。
  \item
        与えられた集合の部分集合は圏を形成する。その圏の射は、最初の集合が二番目の集合の
        部分集合ならば二つの集合を結ぶ矢印として定義される。
        そのような圏での二つの集合の引き戻しは何か？押し出しは何か？
        始対象と終対象は何か？
  \item
        余等化子とは何か推測できるか？
  \item
        終対象を持つ圏において、終対象への引き戻しが積であることを示せ。
  \item
        同様に、始対象からの押し出し（存在するなら）が余積であることを示せ。
\end{enumerate}
